% =====================================================================
%  Chapitre 1 -- Fondements
% =====================================================================
\chapter{Du data warehouse au Lakehouse}\label{ch:fondements}
\soustitre{Principes d'une plateforme de données moderne}

\begin{objectifs}
  \item distinguer les besoins transactionnels et analytiques d'une entreprise ;
  \item expliquer les forces et les limites de l'entrepôt et du lac de données ;
  \item décrire les composants d'un Lakehouse : stockage objet, formats ouverts, catalogue, moteurs ;
  \item expliquer comment un format de table comme Iceberg apporte les transactions et l'historique ;
  \item organiser les données en couches Bronze, Silver et Gold, et les modéliser en étoile ;
  \item situer le traitement par lots et le traitement en continu, et les outils qui exploitent la plateforme.
\end{objectifs}

% =====================================================================
\section{Deux usages des données}
% =====================================================================

Une entreprise de commerce en ligne enregistre chaque jour des milliers de commandes. L'application
qui les reçoit doit répondre en quelques millisecondes à des opérations courtes et nombreuses :
créer une commande, modifier une adresse, enregistrer un paiement. Ce sont des traitements
\textbf{transactionnels}, dits \textbf{OLTP}\index{oltp@OLTP} (\emph{OnLine Transaction Processing}). Ils
s'appuient sur une base relationnelle, comme PostgreSQL ou Oracle, organisée pour garantir la
cohérence de chaque transaction.

La direction, elle, pose d'autres questions : quel est le chiffre d'affaires par région et par
mois ? Quels transporteurs livrent en retard ? Quels produits reçoivent les plus mauvais avis ?
Ces requêtes \textbf{analytiques}, dites \textbf{OLAP}\index{olap@OLAP} (\emph{OnLine Analytical
Processing}), lisent des millions de lignes, sur plusieurs années, pour en tirer quelques
agrégats. Les exécuter sur la base de l'application la ralentirait, et cette base ne garde souvent
que l'état présent des données, pas leur historique.

\begin{table}[H]
\centering\small
\renewcommand{\arraystretch}{1.3}
\begin{tabularx}{\textwidth}{@{}l X X@{}}
\toprule
& \textbf{OLTP (transactionnel)} & \textbf{OLAP (analytique)} \\
\midrule
Question type & Enregistrer la commande n° 4521 & Chiffre d'affaires par mois et par région \\
Volume lu par requête & quelques lignes & des millions de lignes \\
Opérations & insertions et mises à jour fréquentes & lectures et agrégations \\
Historique & état courant & plusieurs années \\
Organisation & par ligne, normalisée & par colonne, souvent dénormalisée \\
Utilisateurs & l'application, des milliers à la fois & analystes, tableaux de bord, modèles \\
\bottomrule
\end{tabularx}
\caption{Les deux usages des données dans une entreprise.}
\label{tab:oltp-olap}
\end{table}

Depuis les années 1990, on sépare donc les deux usages : les données sont copiées régulièrement de
la base transactionnelle vers une plateforme dédiée à l'analyse. Ce chapitre retrace l'évolution de
cette plateforme, de l'entrepôt de données au Lakehouse, puis présente les principes et les outils
mis en œuvre dans le reste de l'ouvrage.

% =====================================================================
\section{L'entrepôt de données}
% =====================================================================

L'\textbf{entrepôt de données}\index{entrepot de donnees@entrepôt de données} (\emph{data warehouse}) est défini par
\citet{inmon2005} comme une collection de données \og orientée sujet, intégrée, historisée et non
volatile \fg{}, destinée à l'aide à la décision. Les données des différentes applications y sont
extraites, nettoyées, rapprochées, puis chargées selon un schéma conçu pour l'analyse : c'est le
processus \textbf{ETL}\index{etl@ETL} (\emph{Extract, Transform, Load}). \citet{kimball2013} a
popularisé la \textbf{modélisation dimensionnelle}, présentée à la section~\ref{sec:dimensionnel},
qui organise ces données en faits et en dimensions.

L'entrepôt offre ce qu'attendent les analystes : du SQL, des performances prévisibles, des
transactions et un schéma contrôlé. Mais il présente trois limites, devenues gênantes avec la
croissance des volumes et la diversité des données :
\begin{itemize}[leftmargin=*,itemsep=2pt]
  \item \textbf{le coût} : stockage et calcul sont liés dans un même système, souvent propriétaire,
        et facturés ensemble ;
  \item \textbf{la rigidité} : seules les données structurées, conformes au schéma, y entrent. Les
        textes, les images, les journaux d'événements restent à l'extérieur ;
  \item \textbf{la fermeture} : les données sont stockées dans un format propre au moteur. Pour
        entraîner un modèle d'apprentissage automatique, il faut les exporter.
\end{itemize}

% =====================================================================
\section{Le lac de données}
% =====================================================================

À la fin des années 2000, l'écosystème Hadoop, puis les services de stockage objet du cloud comme
Amazon S3, rendent le stockage massif très bon marché. L'idée du \textbf{lac de
données}\index{lac de donnees@lac de données} (\emph{data lake}) en découle : stocker \emph{toutes} les données,
dans leur format d'origine (fichiers CSV, JSON, images, journaux), et n'imposer un schéma qu'au
moment de la lecture (\emph{schema-on-read}). Des moteurs de calcul distribué comme Apache
Spark \citep{zaharia2016} traitent ensuite ces fichiers à grande échelle, et les équipes de science
des données y accèdent directement.

Le lac de données a toutefois perdu ce que l'entrepôt garantissait. Un ensemble de fichiers n'est
pas une table : rien n'empêche deux programmes d'écrire en même temps, un traitement interrompu
laisse des fichiers à moitié écrits, une erreur de format passe inaperçue, et il est impossible de
revenir à l'état de la veille. Sans gouvernance, le lac devient un \og marécage \fg{} (\emph{data
swamp}) où personne ne sait plus quelles données sont fiables.

Beaucoup d'entreprises ont donc combiné les deux : un lac pour stocker et préparer, un entrepôt
pour servir les tableaux de bord. Cette \textbf{architecture à deux niveaux} cumule les
inconvénients : les données sont copiées deux fois, avec un décalage, des coûts doublés et un
risque d'incohérence entre les deux systèmes \citep{armbrust2021}.

% =====================================================================
\section{Le Lakehouse}
% =====================================================================

Le \textbf{Lakehouse}\index{lakehouse@Lakehouse} réunit les deux approches. \citet{armbrust2021} le
définissent comme un système de gestion de données fondé sur un stockage à faible coût, en formats
ouverts et directement accessibles, qui offre en plus les fonctions de gestion et de performance
d'un entrepôt : transactions ACID, versions des données, contrôle du schéma, indexation et
optimisation des requêtes. Concrètement, il repose sur quatre briques
(figure~\ref{fig:c1-briques}).

\begin{figure}[H]
\centering
\begin{tikzpicture}[font=\sffamily\small, >=Stealth,
  couche/.style={draw=lhfonce,line width=0.7pt,minimum width=12.6cm,minimum height=1.05cm,align=center},
  et/.style={font=\sffamily\scriptsize,text=lhfonce,align=center}]
\node[couche,fill=lhblue!12] (moteurs) at (0,3.6) {};
\node[et,anchor=west,font=\sffamily\bfseries\small] at (-6.15,3.6) {Moteurs};
\node at (1.2,3.6) {\logo[1.6em]{spark}\quad\logo[1.6em]{trino}\quad\logo[1.6em]{dbt}\quad\logo[1.5em]{superset}\quad\logo[1.5em]{kafka}};
\node[couche,fill=lhvert!25] (cat) at (0,2.3) {};
\node[et,anchor=west,font=\sffamily\bfseries\small] at (-6.15,2.3) {Catalogue};
\node[et] at (1.2,2.3) {associe un nom de table (\texttt{silver.orders}) à ses métadonnées};
\node[couche,fill=lhgold!25] (tab) at (0,1.0) {};
\node[et,anchor=west,font=\sffamily\bfseries\small] at (-6.15,1.0) {Format de table};
\node at (1.2,1.0) {\logo[1.6em]{iceberg}\quad\logo[1.6em]{delta}\quad\textcolor{lhfonce}{\scriptsize transactions, versions, schéma}};
\node[couche,fill=lhsilver!20] (fic) at (0,-0.3) {};
\node[et,anchor=west,font=\sffamily\bfseries\small] at (-6.15,-0.3) {Format de fichier};
\node at (1.2,-0.3) {\logo[1.6em]{parquet}\quad\textcolor{lhfonce}{\scriptsize Parquet : stockage par colonnes, compressé}};
\node[couche,fill=lhgray] (sto) at (0,-1.6) {};
\node[et,anchor=west,font=\sffamily\bfseries\small] at (-6.15,-1.6) {Stockage objet};
\node at (1.2,-1.6) {\logo[1.6em]{rustfs}\quad\textcolor{lhfonce}{\scriptsize compatible S3 : bon marché, illimité}};
\end{tikzpicture}
\caption{Les briques d'un Lakehouse. Chaque couche est ouverte et remplaçable : plusieurs moteurs
lisent les mêmes tables, sans copie.}
\label{fig:c1-briques}
\end{figure}

\subsection{Le stockage objet}

Les données sont rangées dans un \textbf{stockage objet}\index{stockage objet} compatible avec
l'interface S3 d'Amazon : chaque fichier est un objet, identifié par un chemin, dans un espace
appelé \emph{bucket}. Ce stockage est peu coûteux, quasi illimité et indépendant des moteurs de
calcul. Dans cet ouvrage, le service RustFS joue ce rôle sur la machine du lecteur.

\subsection{Un format de fichier par colonnes}

Les données tabulaires sont stockées au format \textbf{Parquet}\index{parquet@Parquet}, un format ouvert
inspiré du système Dremel de Google \citep{melnik2010}. Au lieu de ranger les données ligne par
ligne, comme un fichier CSV, Parquet les range \textbf{colonne par colonne}, par blocs. Une requête
qui n'utilise que trois colonnes sur trente ne lit que ces trois colonnes. Des valeurs voisines
d'une même colonne se ressemblent et se compressent bien. Enfin, chaque bloc conserve le minimum
et le maximum de ses valeurs : un moteur peut ignorer les blocs qui ne contiennent pas la valeur
cherchée.

\subsection{Un format de table ouvert}\label{sec:c1-formats}

Un dossier de fichiers Parquet n'est pas encore une table. Le \textbf{format de
table}\index{format de table} ajoute une couche de métadonnées qui décrit, à chaque instant, quels
fichiers composent la table. Trois formats ouverts dominent : Apache Iceberg, créé chez Netflix,
Delta Lake, créé chez Databricks \citep{armbrust2020}, et Apache Hudi, créé chez Uber. Cet ouvrage
utilise principalement Iceberg, et compare Iceberg et Delta Lake au chapitre~\ref{ch:tp1}.

La figure~\ref{fig:c1-iceberg} montre l'organisation d'une table Iceberg \citep{icebergspec}.
Chaque écriture produit un nouveau fichier de métadonnées, qui décrit un \textbf{instantané}
(\emph{snapshot}\index{snapshot}) de la table : la liste de ses fichiers de données à ce moment-là,
via des fichiers intermédiaires (\emph{manifests}). Les fichiers de données ne sont jamais
modifiés : une mise à jour écrit de nouveaux fichiers, puis un nouvel instantané qui les
référence. Le catalogue pointe vers le fichier de métadonnées courant, et la bascule d'un
instantané au suivant est \textbf{atomique}.

\begin{figure}[H]
\centering
\begin{tikzpicture}[font=\sffamily\footnotesize, >=Stealth,
  b/.style={draw=lhfonce,line width=0.6pt,rounded corners=2pt,align=center,inner sep=4pt,fill=#1},
  et/.style={font=\sffamily\scriptsize,text=lhfonce,align=center}]
\node[b=lhvert!25] (cat) at (0,0) {Catalogue\\\texttt{silver.orders}};
\node[b=lhgold!25] (m1) at (-2.6,-1.7) {\texttt{v1.metadata.json}\\instantané 1};
\node[b=lhgold!25] (m2) at (2.6,-1.7) {\texttt{v2.metadata.json}\\instantanés 1 et 2};
\draw[->,thick,lhfonce] (cat) -- node[et,right=2pt]{version courante} (m2);
\draw[->,lhfonce!50] (m2) -- node[et,above]{précédente} (m1);
\node[b=lhblue!10] (l1) at (-2.6,-3.3) {liste de manifests 1};
\node[b=lhblue!10] (l2) at (2.6,-3.3) {liste de manifests 2};
\draw[->,lhfonce] (m1) -- (l1); \draw[->,lhfonce] (m2) -- (l2);
\node[b=lhblue!5] (ma) at (-1.3,-4.8) {manifest A};
\node[b=lhblue!5] (mb) at (2.6,-4.8) {manifest B};
\draw[->,lhfonce] (l1) -- (ma); \draw[->,lhfonce] (l2) -- (ma); \draw[->,lhfonce] (l2) -- (mb);
\foreach \x/\n in {-3.2/1,-1.3/2,0.6/3} {\node[b=lhsilver!20] (f\n) at (\x,-6.3) {\texttt{f\n.parquet}};}
\node[b=lhsilver!20] (f4) at (2.6,-6.3) {\texttt{f4.parquet}};
\draw[->,lhfonce] (ma) -- (f1); \draw[->,lhfonce] (ma) -- (f2); \draw[->,lhfonce] (ma) -- (f3);
\draw[->,lhfonce] (mb) -- (f4);
\node[et,anchor=west] at (4.1,-1.7) {métadonnées :\\schéma, partitions,\\historique};
\node[et,anchor=west] at (4.1,-4.8) {liste des fichiers\\et statistiques};
\node[et,anchor=west] at (4.1,-6.3) {données (Parquet),\\jamais modifiées};
\end{tikzpicture}
\caption{Une table Iceberg après deux écritures. L'instantané 2 réutilise les fichiers de
l'instantané 1 et ajoute \texttt{f4.parquet} ; l'instantané 1 reste lisible.}
\label{fig:c1-iceberg}
\end{figure}

Cette organisation apporte à de simples fichiers les garanties d'une base de données :
\begin{itemize}[leftmargin=*,itemsep=2pt]
  \item \textbf{des transactions ACID}\index{acid@ACID} : un lecteur voit soit l'ancien instantané, soit
        le nouveau, jamais un état intermédiaire, et deux écritures concurrentes sont arbitrées ;
  \item \textbf{le voyage dans le temps}\index{voyage dans le temps} (\emph{time travel}) : on
        peut relire la table telle qu'elle était à un instantané passé, pour un audit ou pour
        corriger une erreur ;
  \item \textbf{l'évolution du schéma} : on ajoute ou renomme une colonne sans réécrire les données ;
  \item \textbf{le partitionnement caché} : la table est découpée selon une colonne (un mois, un
        vendeur), et une requête filtrée ne lit que les fichiers utiles.
\end{itemize}

En contrepartie, chaque écriture ajoute des fichiers et des métadonnées. Une table alimentée souvent
accumule des \textbf{petits fichiers}\index{petits fichiers}, qu'il faut régulièrement
\textbf{compacter}, et des instantanés anciens, qu'il faut \textbf{expirer}. Les chapitres~\ref{ch:tp3}
et~\ref{ch:tp4} mesurent ce phénomène et le corrigent.

\subsection{Un catalogue et plusieurs moteurs}

Le \textbf{catalogue}\index{catalogue} associe le nom d'une table à son fichier de métadonnées
courant. Tous les moteurs qui passent par le même catalogue voient donc la même version des
données. Cet ouvrage utilise le catalogue REST d'Iceberg. En entreprise, ce rôle est tenu par AWS
Glue, Unity Catalog, Polaris ou Nessie.

C'est la clé de la \textbf{séparation du stockage et du calcul}\index{separation du stockage et du
calcul@séparation du stockage et du
calcul} : les données existent une seule fois, et plusieurs moteurs les utilisent selon leurs
forces. Spark transforme de gros volumes, Trino répond aux requêtes interactives, dbt organise les
transformations SQL, Superset construit les tableaux de bord. Chaque moteur se dimensionne, se met à
jour ou se remplace sans toucher aux données.

% =====================================================================
\section{Organiser les données en couches}\label{sec:c1-medaillon}
% =====================================================================

Dans un Lakehouse, on ne transforme pas les données avant de les charger : on les charge d'abord
telles qu'elles arrivent, puis on les transforme sur place. C'est l'approche
\textbf{ELT}\index{elt@ELT} (\emph{Extract, Load, Transform}), rendue possible par le faible coût du
stockage. Les données traversent alors trois couches successives, selon l'\textbf{architecture en
médaillon}\index{architecture en medaillon@architecture en médaillon} (figure~\ref{fig:c1-medaillon}).

\begin{figure}[H]
\centering
\begin{tikzpicture}[font=\sffamily\footnotesize, >=Stealth,
  et/.style={font=\sffamily\scriptsize,text=lhfonce,align=center,text width=3.3cm}]
\pic at (0,0) {fichier=CSV}; \pic at (0,-1.1) {fichier=\{ \}}; \pic at (0,-2.2) {photo};
\node[et,text width=2cm] at (0,-3.1) {sources};
\pic at (3.3,-1.1) {table=lhbronze}; \node[font=\sffamily\bfseries] at (3.3,-0.25) {Bronze};
\pic at (7.3,-1.1) {table=lhsilver}; \node[font=\sffamily\bfseries] at (7.3,-0.25) {Silver};
\pic at (11.3,-1.1) {table=lhgold};  \node[font=\sffamily\bfseries] at (11.3,-0.25) {Gold};
\pic at (7.3,-3.0) {table=gray};     \node[et,anchor=west,text width=2cm] at (7.8,-3.0) {quarantaine};
\draw[->,very thick,lhfonce] (0.6,-1.1) -- (2.75,-1.1);
\draw[->,very thick,lhfonce] (3.85,-1.1) -- (6.75,-1.1);
\draw[->,very thick,lhfonce] (7.85,-1.1) -- (10.75,-1.1);
\draw[->,thick,lhfonce] (7.3,-1.65) -- (7.3,-2.45);
\node[et] at (3.3,-2.4) {données brutes,\\telles que reçues,\\avec leur origine};
\node[et] at (5.3,-4.1) {};
\node[et] at (11.3,-2.4) {tables métier :\\faits, dimensions,\\indicateurs};
\node[et,text width=2.6cm] at (5.3,-0.55) {typer, contrôler,\\dédoublonner};
\node[et,text width=2.6cm] at (9.3,-0.55) {modéliser,\\agréger};
\end{tikzpicture}
\caption{L'architecture en médaillon. Les données invalides ne sont pas supprimées : elles sont
mises en quarantaine, avec la raison de leur rejet.}
\label{fig:c1-medaillon}
\end{figure}

\begin{description}[leftmargin=1.2em,itemsep=3pt]
\item[Bronze] conserve les données \textbf{telles qu'elles sont arrivées}, avec leur origine et
  leur date de chargement. Si une transformation se révèle fausse, on peut toujours la refaire à
  partir de Bronze, sans redemander les données à la source.
\item[Silver] contient des données \textbf{propres et fiables} : types explicites, valeurs
  contrôlées, doublons écartés, schéma documenté. Les lignes qui ne respectent pas les règles vont
  dans une table de \textbf{quarantaine}\index{quarantaine}, avec la raison de leur rejet, pour
  qu'on puisse les examiner et corriger la source.
\item[Gold] présente des tables \textbf{orientées métier}, prêtes pour l'analyse : tables de faits
  et de dimensions, indicateurs agrégés, données préparées pour un modèle d'apprentissage.
\end{description}

Ce découpage rend chaque étape simple, testable et rejouable. Il sert de fil conducteur aux quatre
travaux pratiques.

% =====================================================================
\section{La modélisation dimensionnelle}\label{sec:dimensionnel}
% =====================================================================

Une base transactionnelle est \textbf{normalisée} : chaque information n'y figure qu'une fois, pour
éviter les incohérences lors des mises à jour. Cette organisation, idéale pour écrire, oblige à de
nombreuses jointures pour analyser. La couche Gold adopte plutôt la \textbf{modélisation
dimensionnelle}\index{modelisation dimensionnelle@modélisation dimensionnelle} de \citet{kimball2013} :
\begin{itemize}[leftmargin=*,itemsep=2pt]
  \item une \textbf{table de faits}\index{table de faits} enregistre des événements mesurables (une
        vente, un avis, une livraison), avec leurs mesures (montant, note, délai) et des clés vers
        les dimensions ;
  \item une \textbf{dimension}\index{dimension} décrit le contexte de ces événements : le client,
        le produit, le vendeur, la date ;
  \item la \textbf{granularité} d'une table de faits précise ce que représente une ligne (une ligne
        de commande, un avis). C'est la première décision de conception.
\end{itemize}

Une table de faits entourée de ses dimensions forme un \textbf{schéma en étoile}\index{schema en
etoile@schéma en
étoile}. Plusieurs tables de faits qui partagent les mêmes dimensions (le même client, la même date)
forment un \textbf{schéma en constellation}\index{schema en constellation@schéma en constellation}, que construit le
chapitre~\ref{ch:tp2} avec l'outil dbt. Les dimensions partagées garantissent que toutes les
analyses parlent du même client et du même calendrier.

% =====================================================================
\section{Au-delà des tables : texte, JSON et images}
% =====================================================================

Une grande part des données d'une entreprise n'est pas tabulaire. Le Lakehouse les accueille sans
les transformer à l'entrée, puis en extrait des informations structurées :
\begin{itemize}[leftmargin=*,itemsep=2pt]
  \item le \textbf{texte libre}, comme les avis clients, se nettoie et se transforme en indicateurs
        ou en variables pour un modèle (chapitre~\ref{ch:tp2}) ;
  \item les données \textbf{semi-structurées}, comme les messages JSON, ont une structure, mais
        imbriquée et variable d'un message à l'autre. On les lit avec un schéma explicite, puis on
        les aplatit en tables (chapitre~\ref{ch:tp2}) ;
  \item les \textbf{images} sont stockées avec leurs octets, puis décrites par leurs dimensions, une
        empreinte qui reconnaît les doublons et un \textbf{vecteur} calculé par un réseau de
        neurones pré-entraîné \citep{howard2019}, qui résume ce qu'elles montrent
        (chapitre~\ref{ch:tp3}).
\end{itemize}

% =====================================================================
\section{Par lots ou en continu}\label{sec:c1-streaming}
% =====================================================================

Les données peuvent être traitées \textbf{par lots}\index{traitement par lots} (\emph{batch}) :
à heure fixe, un traitement lit les données accumulées et produit un résultat complet. C'est simple
et peu coûteux, mais le résultat a toujours le retard de l'intervalle entre deux lots. Le
\textbf{traitement en continu}\index{streaming} (\emph{streaming}) traite au contraire chaque
événement, ou chaque petit groupe d'événements, dès son arrivée.

Le transport des événements est confié à un journal distribué comme \textbf{Apache
Kafka}\index{kafka@Kafka} \citep{kreps2011}. Les événements y sont ajoutés à la fin d'un journal
partitionné, numérotés, et conservés après lecture : plusieurs applications lisent ainsi le même
flux, chacune à son rythme. Spark Structured Streaming \citep{armbrust2018} traite ensuite ce flux
comme une table qui grandit sans fin, par \textbf{micro-lots}.

Le modèle \emph{Dataflow} de \citet{akidau2015} a formalisé la principale difficulté du traitement
en continu : les événements n'arrivent pas toujours dans l'ordre où ils se sont produits. Il faut
distinguer l'\textbf{heure de l'événement}\index{heure de l'evenement@heure de l'événement} de l'heure de son
traitement, et décider combien de temps attendre les retardataires, avec un
\textbf{watermark}\index{watermark}. Ce choix règle un compromis entre la fraîcheur et l'exactitude
des résultats \citep{akidau2018}. Le chapitre~\ref{ch:tp4} le met en œuvre, et compare les deux
approches à partir de mesures.

% =====================================================================
\section{Exploiter la plateforme}
% =====================================================================

Une plateforme de données ne se limite pas au stockage et aux transformations. Cinq fonctions
complètent l'architecture, et chacune fait l'objet d'une partie des travaux pratiques :
\begin{itemize}[leftmargin=*,itemsep=2pt]
  \item l'\textbf{orchestration}\index{orchestration} enchaîne les traitements, selon un calendrier
        ou un événement, les relance en cas d'erreur passagère et en garde l'historique (Airflow) ;
  \item le \textbf{SQL interactif} et la \textbf{fédération} permettent aux analystes d'interroger
        le Lakehouse, et d'autres bases, en une seule requête (Trino) ;
  \item la \textbf{visualisation} met les indicateurs à la disposition des équipes métier (Superset,
        Streamlit) ;
  \item la \textbf{supervision} mesure l'état des services et déclenche des alertes (Prometheus,
        Grafana) ;
  \item l'\textbf{observabilité des données}\index{lineage} retrace l'origine de chaque table, son
        \emph{lineage}, à partir des exécutions réelles (OpenLineage, Marquez).
\end{itemize}

% =====================================================================
\section{Les outils de cet ouvrage}
% =====================================================================

Le tableau~\ref{tab:c1-outils} associe chaque fonction aux logiciels libres utilisés dans les
travaux pratiques, et à des équivalents courants en entreprise. Les notions s'y transposent
directement.

\begin{table}[H]
\centering\small
\renewcommand{\arraystretch}{1.28}
\begin{tabularx}{\textwidth}{@{}>{\raggedright\arraybackslash}p{3.4cm} >{\raggedright\arraybackslash}X >{\raggedright\arraybackslash}X c@{}}
\toprule
\textbf{Fonction} & \textbf{Dans cet ouvrage} & \textbf{En entreprise} & \textbf{Chapitre} \\
\midrule
Stockage objet & RustFS (compatible S3) & Amazon S3, Azure Data Lake Storage & \ref{ch:environnement} \\
Format de table & Apache Iceberg, Delta Lake & Iceberg, Delta Lake, Apache Hudi & \ref{ch:tp1} \\
Catalogue & Iceberg REST & AWS Glue, Unity Catalog, Polaris & \ref{ch:tp1} \\
Calcul & Apache Spark & Databricks, Amazon EMR, Snowflake & \ref{ch:tp1} \\
Transformations SQL & dbt & dbt Cloud & \ref{ch:tp2} \\
Orchestration & Apache Airflow & Airflow géré, Dagster, Prefect & \ref{ch:tp3} \\
SQL interactif & Trino & Amazon Athena, Starburst & \ref{ch:tp3} \\
Tableaux de bord & Apache Superset, Streamlit & Power BI, Tableau & \ref{ch:tp1}, \ref{ch:tp3} \\
Supervision & Prometheus, Grafana & Datadog, CloudWatch & \ref{ch:tp3} \\
Lineage & OpenLineage, Marquez & DataHub, Unity Catalog & \ref{ch:tp2}, \ref{ch:tp3} \\
Flux d'événements & Apache Kafka & Confluent Cloud, Amazon MSK & \ref{ch:tp4} \\
Traitement en continu & Spark Structured Streaming & Apache Flink, Kafka Streams & \ref{ch:tp4} \\
\bottomrule
\end{tabularx}
\caption{Les fonctions d'une plateforme de données, et les outils qui les assurent.}
\label{tab:c1-outils}
\end{table}

% =====================================================================
\section{Le fil conducteur : Olist}
% =====================================================================

Tous les travaux pratiques portent sur la même entreprise : \textbf{Olist}\index{olist@Olist}, une place
de marché brésilienne qui relie de petits vendeurs aux grands sites de commerce en ligne. Olist a
publié un jeu de données réel et anonymisé d'environ 100\,000 commandes passées entre 2016 et 2018
\citep{olist2018} : commandes, articles, paiements, clients, vendeurs, produits, et avis des clients
en portugais. Les chapitres l'enrichissent progressivement de données reconstituées à partir de
ces commandes (messages de suivi des colis, flux d'événements) et de photos de produits issues du
jeu \emph{Amazon Berkeley Objects} \citep{collins2022}. L'annexe~\ref{ann:donnees} décrit ces
données et leurs licences.

Travailler sur une seule entreprise a un intérêt pédagogique : chaque nouveau besoin s'appuie sur
ce qui a été construit avant, comme dans un vrai projet. Les tables Silver du chapitre~\ref{ch:tp1}
alimentent la constellation du chapitre~\ref{ch:tp2}, que le pipeline automatisé du
chapitre~\ref{ch:tp3} reconstruit chaque jour, et que le flux du chapitre~\ref{ch:tp4} complète en
continu.

\begin{aretenir}
  \item Les usages transactionnels (OLTP) et analytiques (OLAP) demandent des systèmes différents.
  \item L'entrepôt de données est fiable mais coûteux et fermé. Le lac de données est ouvert et bon
        marché, mais sans garanties.
  \item Le Lakehouse ajoute à un stockage objet en formats ouverts (Parquet) une couche de
        métadonnées (Iceberg, Delta Lake) qui apporte transactions, versions et contrôle du schéma.
  \item Un catalogue partagé permet à plusieurs moteurs d'utiliser les mêmes tables, sans copie.
  \item Les données traversent trois couches : Bronze (brut), Silver (fiable), Gold (métier), et la
        couche Gold adopte la modélisation dimensionnelle.
  \item Le traitement en continu doit gérer le désordre des événements : le watermark règle le
        compromis entre fraîcheur et exactitude.
\end{aretenir}

\begin{plusloin}
L'article fondateur de \citet{armbrust2021} présente les motivations du Lakehouse, et celui de
\citet{armbrust2020} le fonctionnement interne de Delta Lake. La spécification d'Iceberg
\citep{icebergspec} détaille les métadonnées et les garanties de concurrence. Pour la modélisation
dimensionnelle, l'ouvrage de \citet{kimball2013} reste la référence. \citet{kleppmann2017} offre
une vue d'ensemble des systèmes de données, des bases transactionnelles aux flux, et
\citet{reis2022} une introduction au métier d'ingénieur de données.
\end{plusloin}
